\documentclass[11pt]{article}
\usepackage[margin=25mm]{geometry}
\usepackage{fontspec}
\setmainfont{Latin Modern Roman}
\usepackage{amsmath,amssymb,mathtools}
\usepackage{unicode-math}
\setmathfont{Latin Modern Math}
\usepackage{xurl,hyperref}
\hypersetup{colorlinks=true,urlcolor=blue}
\setcounter{secnumdepth}{0}
\setlength{\parindent}{0pt}
\setlength{\parskip}{5pt}
\setlength{\emergencystretch}{3em}
\newcommand{\sref}[2]{\hyperlink{source:#1:#2}{[S#2]}}
\newcommand{\cataloguescope}[1]{\paragraph{Recorded target.}#1}
\newcommand{\R}{\mathbb R}
\newcommand{\C}{\mathbb C}
\newcommand{\N}{\mathbb N}
\newcommand{\Z}{\mathbb Z}
\newcommand{\Q}{\mathbb Q}
\newcommand{\F}{\mathbb F}
\newcommand{\eps}{\varepsilon}
\begin{document}
\section{Phase Diagram of the Interacting Bose Gas}\label{statement}

\subsection{Definitions and mathematical statement}

\paragraph{Shared notation.}$\mathbb N=\{1,2,\ldots\}$; the standard number fields have their usual meanings. All additional conventions are stated locally.


\paragraph{Definitions and conventions.}
In a grand-canonical Bose gas, chemical potential $\mu$ controls particle density at fixed temperature and interaction. Correlation clustering means that separated observables become independent; exponential and power-law clustering specify the distance dependence. Bose--Einstein condensation can be expressed by a macroscopic eigenvalue of the one-particle density matrix, with off-diagonal long-range order in a homogeneous setting. The proposed phases are distinguished by these large-distance properties in the thermodynamic limit.

\paragraph{Target retained from the canonical source catalogue.}
\paragraph{Catalogue wording retained for traceability.}
Establish the conjectured phase diagram for the three-dimensional interacting Bose gas: below the critical chemical potential the correlations should cluster exponentially, at criticality by a power law, and above it the system should exhibit Bose–Einstein condensation with no clustering.

\paragraph{Source scope and qualification.}
The interaction and thermodynamic parameters are source-specific. The nonclustering assertion refers to the condensate correlation channel, not a claim that every connected observable correlation ceases to decay.
\subsection{Short English statement}
Prove the predicted three-dimensional Bose-gas phases: exponential decay below criticality, power laws at criticality, and condensation above it.
\subsection{Research attempt}
\paragraph{Model, date and investigated scope.}
OpenAI GPT-6 Astra (configured model \texttt{gpt-6-astra}), ultra reasoning; generated 11 October 2026. Derives the massive and massless three-dimensional Gaussian kernels and separates condensate correlations from connected clustering; all interacting Bose-gas phases remain outside the proof.

\paragraph{Sources and literature check.}
The complete supplied definition, including source qualifications and the dated literature assessment, was reviewed. The calculation below is an elementary restricted attempt motivated by [S2]; it does not claim novelty or a fresh literature-status determination.

\paragraph{Context and investigated scope.} The complete supplied definition, including source qualifications and the dated literature assessment, was reviewed. The calculation below is an elementary restricted attempt motivated by [S2]; it does not claim novelty or a fresh literature-status determination.

\paragraph{A Gaussian test of the proposed three regimes.}
As an auxiliary static model in three dimensions, consider the covariance of $-\Delta+m^2$, where $m\geq0$. For $r=|x|>0$, a decaying radial Green function satisfies
\[
 -G_m''(r)-\frac2rG_m'(r)+m^2G_m(r)=0,
 \qquad G_m(r)=\frac{e^{-mr}}{4\pi r}.
\]
Indeed, putting $v=rG_m$ reduces the equation to $v''=m^2v$. Decay removes the growing solution when $m>0$, and the delta-source normalization follows from
$\lim_{r\downarrow0}-4\pi r^2G_m'(r)=1$.
For $m=0$, the solution decaying at infinity is $1/(4\pi r)$, with the same normalization. Thus the massive covariance has exponential decay with correlation length $m^{-1}$; removing the mass leaves a power law. This elementary calculation explains a possible mechanism for the first two regimes, without controlling an interacting measure.

\paragraph{Which correlation does not cluster?}
Let $\eta$ be a centered complex Gaussian generalized field with covariance $G_m$, and formally write $\psi=z+\eta$ for a constant $z\in\mathbb C$. For distinct points, or after smearing, its two-point kernel is
\[
 \mathbb E\bigl[\overline{\psi(x)}\psi(y)\bigr]
   =|z|^2+G_m(x-y).
\]
The connected kernel, obtained by subtracting the product of means, is still $G_m$. A nonzero constant contribution therefore coexists with clustering of the connected kernel. One can also average the phase of $z$ uniformly: the mean then vanishes while the constant two-point contribution survives. This tests the definition's qualification concerning the condensate channel, rather than asserting failure of every connected correlation.

\paragraph{Failure of the Gaussian continuation.}
In the classical quadratic approximation with kinetic operator $-\Delta$, chemical potential enters as $m^2=-\mu$. On a finite periodic box, its constant mode has quadratic coefficient $-\mu$. For $\mu>0$ the Gaussian integral diverges in that direction, and at $\mu=0$ the unconstrained finite-volume zero-mode integral also diverges. Consequently the calculation cannot simply be continued to produce the condensed phase. Adding a prescribed $z$ is not a derivation of a condensate from interactions. An interaction-dependent critical chemical potential, the thermodynamic construction, critical exponents and spontaneous condensation all remain unproved. In particular, the kernel above is a static Gaussian surrogate, not the full quantum Bose-gas correlation function.

\paragraph{Verification and adversarial audit.}
Checked the radial Green equation and delta normalization, the distinction between centered and phase-averaged correlations, and the finite-volume zero-mode divergence. Model self-audit only; no independent human proof checking. Prompts A and B were applied to this canonical record and its qualified source scope. Checked the radial Green equation and delta normalization, the distinction between centered and phase-averaged correlations, and the finite-volume zero-mode divergence. Retained unresolved status and zero original-target progress; Construct and control the interacting thermodynamic measure near its critical chemical potential; the Gaussian integral already fails on the positive-chemical-potential constant mode.

Coordinator checked the radial Green-function equation and delta normalization, the connected-versus-condensate correlation distinction and the finite-box Gaussian zero-mode obstruction. This additional check was model-assisted, not independent human verification.

\paragraph{Tools, process and contributions.}
Python 3 was used for file preparation and structural checks; no numerical experiment or theorem prover was used. Contribution prompt A was followed by prompt B. GPT-6 Astra with ultra reasoning drafted the research and performed its stated self-audit; any additional coordinator contribution is identified in the verification and audit. A GPT-6-based Codex coordinator assembled the manuscript and checked repository integration; its exact variant and reasoning configuration were not exposed. No human mathematical contribution, independent human proof check or formal proof certification is supplied. Complete derivations and the consulted source references are included here; private reasoning transcripts are not published.

\paragraph{Final status and first remaining gap.}
\textbf{UNRESOLVED}. Construct and control the interacting thermodynamic measure near its critical chemical potential; the Gaussian integral already fails on the positive-chemical-potential constant mode.

\section{Completion Estimate}
\noindent\textbf{COMPLETION: 0\%}
The calculation is an auxiliary or restricted attempt and establishes no newly verified portion of the full source target. Construct and control the interacting thermodynamic measure near its critical chemical potential; the Gaussian integral already fails on the positive-chemical-potential constant mode.

\subsection{Sources}
\begin{itemize}
\item[\hypertarget{source:statement:1}{S1}] ULAM.AI and the UnsolvedMath Contributors, UnsolvedMath, catalogue entry. Catalogue curation under CC BY 4.0; definitions and English exposition added locally.
\url{https://huggingface.co/datasets/ulamai/UnsolvedMath}.
\item[\hypertarget{source:statement:2}{S2}] The Renormalization Group (2016), Oberwolfach Reports, DOI 10.4171/owr/2016/26, pp. 1453--1515. Primary report accessed 2026-10-11; relevant PDF pages 34.
\url{https://doi.org/10.4171/owr/2016/26}.
\end{itemize}
\label{end:statement}
\end{document}
